\subsection{Implementation and finite evidence for profile discovery}
\label{sec:profile-evidence}

The general component-budget search is implemented in
\path{normal_weight_clusters.py}. It enumerates connected supports,
constructs exact local affine profiles, retains admissible positive even
increments, and reuses those profiles in disjoint component products.
It also implements the dependent-baseline homogeneous exception.
Projected rows are deduplicated, and a same-sign test rejects equations
that cannot have strictly positive increments.  All source conditions
are checked.  The component budget is capped at $\min(r,k)$, including
when the supplied integer $r$ is very large.  Global work, support,
linear-system, and product limits remain effective through final
certificate replay; a limit reports incomplete search.

The supplied vertex gate is an unchanged copy of the incoming
rank-and-boundary checker.  It verifies a primitive supported normal
ray, Euler characteristic one, and essential boundary through exact
linear algebra and boundary cohomology.  Positive results are replayed
by the separate checker.  Exhaustion concerns only the stated source
and parameter class and yields no unrestricted knot verdict.  The
overlay leaves the maintained recognizer and its defaults unchanged.

\subsubsection{Complete radius-two candidate generation}

The paired experiment compares the new profile search with the incoming
enumerator over all supports of size at most two.  Each timed call
includes source and cocycle preparation, complete candidate generation,
local normal reconstruction, and actual-cancellation filtering; neither
arm includes a topological vertex-disc gate.  Candidate sets are checked
for exact equality outside the timed blocks.  Each source receives one
excluded warmup per arm and five measured rounds, with arm order
shuffled using a recorded seed.  Table~\ref{tab:profile-radius-two}
reports seconds and ratios of the two arm medians; these are not medians
of per-round ratios.  Full samples, observed ranges, counters, and hashes
are in \texttt{artifacts/benchmark.json}.

\begin{table}[htbp]
\centering
\small
\begin{tabular}{lrrrr}
\hline
Source & $t$ & All supports (s) & Profiles (s) & Ratio\\
\hline
Layered obstruction & 8 & 1.189400 & 0.026278 & $45.26$\\
Layered obstruction & 16 & 7.213136 & 0.043151 & $167.16$\\
Layered obstruction & 24 & 25.963222 & 0.060250 & $430.93$\\
Layered obstruction & 40 & 122.938084 & 0.096843 & $1269.46$\\
Disconnected-support source & 15 & 4.006094 & 0.041128 & $97.40$\\
\hline
\end{tabular}
\caption{Complete radius-two candidate generation, five measured
rounds.  The ratio is the all-support median divided by the profile
median.  Exact resulting candidate sets agree in every case.}
\label{tab:profile-radius-two}
\end{table}

The counts make the change in workload concrete.  At $t=40$, the
incoming arm examines $904$ weight spaces and $156\,927$ small systems;
the profile arm examines $173$ connected supports and $4\,051$ systems.
Only two local profiles survive, producing three assembled candidates
including the coherent candidate.  For the disconnected-support source,
the corresponding counts are $154$ spaces and $14\,942$ systems versus
$78$ connected supports and $1\,420$ systems.  Three reusable local
profiles produce five assembled candidates.  Maximum tetrahedron-edge
incidence is $d=6$ in the smallest layered source and $d=8$ in the other
sources.

These gains combine component reuse, smaller connected-support
enumeration, projected-row deduplication, and the stricter local
positive-even-increment filters.  The old baseline first searches
homogeneous weight spaces and filters the actual cancellation support
after reconstruction.  Thus the table does not isolate connectedness
as the sole cause of the speedup.  The source family is deliberately
structured, the parameter radius is small, and these finite timings
establish neither a corpus-wide speedup nor an asymptotic growth law.

\subsubsection{A retained whole-query regression}

The existing incoming one-edge specialization is a stronger practical
baseline for positive radius-one queries.  A second paired experiment
therefore compares complete successful queries, including internal
positive verification and replay.  Both methods find the same certified
disc on these sources.  Here there is one excluded warmup per arm and
seven measured rounds, again with shuffled arm order.  The incoming
method uses its native component observer, whereas the new method uses
the polynomial rank gate.  Their end-to-end costs consequently include
different verification mechanisms, an intentional part of this
integration comparison.

\begin{table}[htbp]
\centering
\small
\begin{tabular}{rrrr}
\hline
$t$ & Incoming one-edge (s) & Profiles with rank gate (s) & Slowdown\\
\hline
8 & 0.002818 & 0.005381 & $1.91$\\
16 & 0.004930 & 0.014919 & $3.03$\\
24 & 0.007299 & 0.034642 & $4.75$\\
40 & 0.012197 & 0.128793 & $10.56$\\
\hline
\end{tabular}
\caption{Complete positive radius-one queries, seven measured rounds.
The slowdown is the new median divided by the incoming median, not the
median of paired sample ratios.}
\label{tab:profile-radius-one}
\end{table}

The larger-radius search improvement therefore does not justify
replacing the existing one-edge specialization.  A practical portfolio
should retain it, while adding component profiles for larger budgets.
The polynomial bit-complexity guarantee of the rank gate also does not
imply lower constants on these small positive instances.  All measured
regressions are retained in the artifact.

\subsubsection{Validation and evidential limits}

All eleven focused tests pass.  They compare exhaustive candidate sets
on actual sources; recover the disconnected witness through the
independent incoming weight-space enumerator; exercise the homogeneous
exception; compare connected-set enumeration with brute subsets; check
global resource limits, external cancellation, and input immutability;
replay certificates with discovery routines disabled; reject forged
certificates; and exercise exact arithmetic with an $8193$-bit right
hand side.  That last test validates the local arithmetic path, not a
complete topological search at that coordinate size.

The retained fifteen-tetrahedron source includes its full face
pairings, construction history, independently checked Pachner moves,
and positive vertex-disc certificates.  Its independently generated
dense one-edge inventory proves the finite statement
$\kappa_V=2$ as described in
Proposition~\ref{prop:cluster-fixture-minimum}.  The separate full census
contains $821$ standard vertex vectors agreed upon by Regina's direct
tree and direct double-description algorithms; the three meridians all
have two cancellation components.  Each positive is independently
certified.  Completeness of this additional census still relies on
those external enumeration algorithms.  No generic independently
replayable certificate for unrestricted negative assertions is supplied.

The new module, tests, measurement driver, finite-audit scripts, and
theorem fragment form an additive research artifact.  They establish a
proved search guarantee for an explicit source class and measured gains
for complete radius-two candidate generation, together with a concrete
integration regression.  A global quasipolynomial recognizer still
requires the structural and source-coverage results stated elsewhere in
this article.
